\documentclass[10pt,a4paper]{amsart}
\usepackage[margin=19mm]{geometry}
\usepackage{amsmath,amssymb,mathtools}
\usepackage[scr=rsfs,cal=euler]{mathalfa}
\usepackage{xfrac}
\usepackage[unicode,hidelinks,bookmarks=false]{hyperref}
\newcommand{\ie}{i.e.\ }
\newcommand{\cf}{cf.\ }
\newcommand{\resp}{respectively\ }
\newcommand{\Subsection}{\subsection}
\providecommand{\nobreakdash}{\nobreak\hbox{-}}
\setlength{\emergencystretch}{2em}
\multlinegap0pt
\usepackage{bm}
\usepackage{bbm}
\usepackage{dsfont}
\usepackage{xspace}
\usepackage{cancel}
\usepackage{mathtools}
\usepackage{amsthm}
\makeatletter
\@ifundefined{thm}{\newtheorem{thm}{Thm}}{}
\@ifundefined{lem}{\newtheorem{lem}[thm]{Lemma}}{}
\makeatother
\newcommand{\PiBC}{\Pi^{\msfBC}}
\newcommand{\bn}{\mathfrak{B}_n}
\newcommand{\dfct}{\textsc{dfct}}
\newcommand{\inv}{\textsc{inv}}
\newcommand{\mfs}[1]{\mathfrak{S}_{#1}}
\newcommand{\msfBC}{\mathsf{BC}}
\newcommand{\nTksp}{\negthickspace\negthickspace}
\newcommand{\ol}[1]{\overline{#1}}
\newcommand{\sumsb}[1]{\sum_{\substack{#1}}}  %%%for multi-line sums
\newcommand{\type}{\mathrm{type}}
\title[A generalization of Deodhar's defect statistic for Iwahori--]{A generalization of Deodhar's defect statistic for Iwahori--Hecke algebras of type $BC$}
\author{Gavin Hobbs, Tommy Parisi, Mark Skandera,, Jiayuan Wang}
\date{Source: arXiv:2506.20099v1 (CC BY 4.0)}
\begin{document}
\maketitle

\noindent\emph{Source and licence.} A generalization of Deodhar's defect statistic for Iwahori--Hecke algebras of type $BC$. Version arXiv:2506.20099v1 (CC BY 4.0), distributed under the Creative Commons Attribution 4.0 International licence, as recorded in the arXiv licence metadata for this exact version and shown on the abstract page https://arxiv.org/abs/2506.20099v1. Full rights notice: licenses/arxiv-2506.20099-CC-BY-4.0.txt.\par\smallskip

\noindent\textit{Editorial scope.} Counted unit: the Lem whose source label is \texttt{l:collapse} in the article (source0.tex lines 2113--2123, source label \texttt{l:collapse}), delivered together with the article's own complete original proof of it (source0.tex lines 2124--2190, one \texttt{proof} environment of 67 lines). No mathematics is written, repaired or re-derived: statement and proof are the article's own text line for line. Required context retained uncounted: none. Every definition, display and label of those lines is retained unchanged. Omitted from this package and not redistributed: 1--2112, 2191--2337. Nothing inside a retained excerpt is omitted or altered: every retained body is delivered exactly as the article's lines are written, so no omission receipt of any kind accompanies this package. Citations: the retained excerpts carry no citation command at all, so no bibliography entry of the article is transcribed and no reference is redistributed. Reference adaptation: none was needed and none was made; every reference inside the delivered unit resolves inside it, so no unresolved reference or citation is produced. Printed slips kept literal: no printed word, symbol, hypothesis or number of the counted statement or of its proof was altered. Layout and class: the article's own class is \texttt{amsart} with packages including amsmath, amsthm, graphicx, amssymb, enumerate, tikz, hyperref, cleveref, tabularx, mathrsfs; the wrapper is the lane's pinned \texttt{amsart} editorial wrapper at 10pt on A4, which restates the theorem-like environments the retained text needs and adds the article's own macro definitions that the retained text needs (\texttt{\textbackslash PiBC}, \texttt{\textbackslash bn}, \texttt{\textbackslash dfct}, \texttt{\textbackslash inv}, \texttt{\textbackslash mfs}, \texttt{\textbackslash msfBC}, \texttt{\textbackslash nTksp}, \texttt{\textbackslash ol}, \texttt{\textbackslash sumsb}, \texttt{\textbackslash type}), with extra wrapper packages (bm, bbm, dsfont, xspace, cancel, mathtools). The delivered numbering is the unit's own. Assets: no retained span contains a figure environment, a graphics-inclusion command, a PDF-page inclusion, a table or a TikZ picture; no image file is copied into this package, the package declares no asset dependency and no image file is redistributed with it. Licence: version arXiv:2506.20099v1 of this article is distributed under the Creative Commons Attribution 4.0 International licence, as recorded in the arXiv licence metadata for this exact version and shown on the abstract page https://arxiv.org/abs/2506.20099v1. A full rights notice is delivered at licenses/arxiv-2506.20099-CC-BY-4.0.txt. Characters: every retained excerpt is pure ASCII, so the delivered engine, XeLaTeX with its default Latin Modern text font, reports no missing character.
\begin{lem}\label{l:collapse}
Fix a generalized star network $G$ in which $m$ multiplicity-$1$ edges are incident upon a pair $\{x,y\}$ of vertices and the $m$ reflections of these edges are incident upon the reflections $\{\ol x, \ol y\}$ of $x$ and $y$. Assume that at most one of the vertices $\{ x, y \}$ equals its reflection. Replace the edges incident upon $x$ and $y$
by a single multiplicity-$m$ edge,
replace their reflections by another such edge, and call the resulting generalized star network $F$.
Then we have
\begin{equation}\label{eq:mto1}
        \sum_{\pi \in \PiBC(G)} \nTksp q^{\dfct^{\msfBC}(\pi)} T_{\type(\pi)} = 
    [m]_q! \nTksp
    \sum_{\pi \in \PiBC(F)} \nTksp q^{\dfct^{\msfBC}(\pi)} T_{\type(\pi)}.
\end{equation}
\end{lem}

\begin{proof}
Assume that $x$ and $y$ lie on or above the horizontal line of symmetry in $G$. Let $\mathcal E = \{\varepsilon_1,\dotsc, \varepsilon_m\}$ be the set of $m$ edges from $x$ to $y$, labeled from bottom to top, and let $\ol{\mathcal E}$ be the set of reflections of these.  

Consider a path family $\pi \in \PiBC(G)$ with the property that
edges $\varepsilon_1,\dotsc,\varepsilon_m$ belong to
paths $\pi_{i_1},\dotsc,\pi_{i_m}$, 
respectively, with $i_1 < \cdots < i_m$.
Define 
$\Phi(\pi,\mathcal E)$ to be the set of all path families in $\PiBC(G)$ which agree with $\pi$ on every edge of $G$, except possibly for those edges belonging to $\mathcal E \cup \ol{\mathcal E}$.
Since each path family $\tau \in \Phi(\pi,\mathcal E)$ 
can differ from $\pi$
only in that 
edges $\epsilon_1,\dotsc,\epsilon_m$ belong to 
paths $\tau_{u_1},\dotsc,\tau_{u_m}$, respectively,
for some permutation $u_1 \cdots u_m$ of $\{i_1,\dotsc,i_m\}$, we have that
$\type(\tau) = \type(\pi)$ and 
$|\Phi(\pi,\mathcal E)| = m!$.
  

We claim that each path family $\tau \in \Phi(\pi,\mathcal E)$ has defects which agree with those of $\pi$ except possibly at the vertex $y$, where $\pi$ has no defect:
   \begin{equation}\label{eq:taupidefect}
    \dfct^{\msfBC}(\tau) = \dfct^{\msfBC}(\pi) + \# \{ \text{defects of $\tau$ occurring at vertex $y$}\}.
\end{equation}
To see this, observe that our choice of $\pi$ implies that each pair of paths among $\pi_{i_1},\dotsc,\pi_{i_m}$ crosses an even number of times before meeting at $y$.  Also the source-to-$x$ subpaths of $\pi_{i_1},\dotsc,\pi_{i_m}$ and $\tau_{i_1},\dotsc,\tau_{i_m}$ are identical. Now any vertex $z$ to the right of $y$
belongs to
paths $\tau_k$, $\tau_l$ if and only if it belongs to $\pi_k, \pi_l$.
Furthermore, if $(\pi_k, \pi_l)$ and $(\tau_k,\tau_l)$ appear in the same relative order between $x$ and $y$, then the number of crossings of $(\pi_k, \pi_l)$ preceding $z$ equals the number of crossings of $(\tau_k, \tau_l)$ preceding $z$. Otherwise these numbers of crossings differ by $2$, with one pair crossing at $x$ and at $y$ while the other pair crosses at neither of these vertices.
In either case, the meeting of $(\pi_k, \pi_l)$ at $z$ is defective if and only if the meeting of $(\tau_k, \tau_l)$ at $z$ is defective.


Consider therefore the defects occurring at vertex $y$. By our choice of $i_1,\dotsc,i_m$, %Then 
the number of defects of $\tau$ occurring at vertex $y$ equals $\inv(u_1 \cdots u_m)$. Thus by (\ref{eq:taupidefect}) we have 

 \begin{equation}\label{eq:mqfactorial}
     \sum_{\tau \in \Phi(\pi,\mathcal E)} q^{\dfct^{\msfBC}(\tau)} = q^{\dfct^{\msfBC}(\pi)}\nTksp\nTksp
     \sum_{u \in \mfs{\{i_1,\dotsc,i_m\}}} \nTksp \nTksp q^{\inv(u)}
     %q^{\dfct(\pi)} 
     = 
     [m]_q! q^{\dfct^{\msfBC}(\pi)},
 \end{equation}
 and the coefficient of $T_w$ on left-hand-side of (\ref{eq:mto1}) equals
 \begin{equation}\label{eq:nodefectaty}
     %\sum_{w \in \bn} T_w \nTksp 
     \sum_{\tau \in \PiBC_w(G)}\nTksp q^{\dfct^{\msfBC}(\tau)} =
     %\sum_{w \in \bn} T_w
     \nTksp\sumsb{\pi \in \PiBC_w(G)\\
     \text{having no }\\\text{defect at } y}
     %q^{\dfct(\pi)} 
     \sum_{\tau \in \Phi(\pi,\mathcal E)} q^{\dfct^{\msfBC}(\tau)}
     = [m]_q! 
     %\sum_{w \in \bn} T_w \nTksp 
     \nTksp\sumsb{\pi \in \PiBC_w(G)\\ 
     \text{having no }\\\text{defect at } y}\nTksp 
     %[m]_q! 
     q^{\dfct^{\msfBC}(\pi)}.
 \end{equation}
 Now create network $F$ from $G$ by deleting edges $\varepsilon_2,\dotsc,\varepsilon_n$, $\ol{\varepsilon_2},\dotsc,\ol{\varepsilon_n}$, and by assigning multiplicity $m$ to $\varepsilon_1$ and $\ol{\varepsilon_1}$. 
 Similarly, create path family $\pi' \in \PiBC(F)$ from $\pi$ by replacing edges $\varepsilon_2,\dotsc,\varepsilon_m$ in $\pi_{i_2},\dotsc,\pi_{i_m}$
 with $\varepsilon_1$, and edges $\ol{\varepsilon_2},\dotsc,\ol{\varepsilon_m}$ 
 in $\pi_{\ol{i_1}},\dotsc,\pi_{\ol{i_m}}$
 with $\ol{\varepsilon_1}$. Specifically we have $\pi' \in \PiBC_{\type(\pi),\dfct(\pi)}(F)$. Moreover, it is easy to see that for each pair $(w,d) \in \bn \times \mathbb N$, this procedure defines a bijective correspondence between the sets
 $\{ \pi \in \PiBC_{w,d}(G) \,|\, \pi \text{ has no defect at } y \}$ and $\PiBC_{w,d}(F)$. Thus we may rewrite the final sum in (\ref{eq:nodefectaty}) as
\begin{equation*}
    [m]_q! \nTksp\sum_{\pi \in \PiBC_w(F)}\nTksp q^{\dfct^{\msfBC}(\pi)},
\end{equation*}
which is the coefficient of $T_w$ on the right-hand-side of (\ref{eq:mto1}).
\end{proof}

\end{document}
